\section{Ôn tập chương 9 - Đạo hàm}
\Opensolutionfile{ans}[ans/ans-OTC9-TN]
\caulc
\begin{ex}%Câu 1.%[1D7N1-1]
	Cho hàm số $y=f(x)$ xác định trên $\mathbb{R}$ thỏa mãn $\lim\limits_{x\to 4}\dfrac{f(x)-f(4)}{x-4}=5$. Kết quả đúng là
	\choice
	{$f'(2)=3$}
	{$f'(x)=2$}
	{$f'(x)=5$}
	{\True $f'(4)=5$}
	\loigiai{
		Theo định nghĩa đạo hàm của hàm số tại một điểm ta có
		$\lim\limits_{x\to 4}\dfrac{f(x)-f(4)}{x-4}=5=f'(4)$.}
\end{ex}
\begin{ex}%Câu 2.%[1D7N2-1]
	Cho hàm số $y=\dfrac{4}{x-1}$. Khi đó $y'(-1)$ bằng
	\choice
	{\True $-1$}
	{$-2$}
	{$2$}
	{$1$}
	\loigiai{
		Ta có $y'=-\dfrac{4}{(x-1)^2}\Rightarrow y'(-1)=-1$.}
\end{ex}
\begin{ex}%Câu 3.%[1D7H2-1]
	Cho hàm số $f(x)=\sqrt{x^2+3}$. Tính giá trị của biểu thức $S=f(1)+4f'(1)$. 
	\choice
	{\True $S=4$}
	{$S=2$}
	{$S=6$}
	{$S=8$}
	\loigiai{
		Ta có: $f(x)=\sqrt{x^2+3}\Rightarrow f'(x)=\dfrac{x}{\sqrt{x^2+3}}$.\\
		Vậy $S=f(1)+4f'(1)=4$.}
\end{ex}
\begin{ex}%Câu 4.%[1D7N2-3]
	Tiếp tuyến của đồ thị hàm số $y=\dfrac{x+1}{2x-3}$ tại điểm có hoành độ $x_0=-1$ có hệ số góc bằng
	\choice
	{$5$}
	{\True $-\dfrac{1}{5}$}
	{$-5$}
	{$\dfrac{1}{5}$}
	\loigiai{
		Tập xác định $\mathscr{D}=\mathbb{R}\setminus\left\{\dfrac{3}{2}\right\}$.\\
		Ta có $f'(x)=\dfrac{-5}{(2x-3)^2}$.\\
		Hệ số góc của tiếp tuyến của đồ thị hàm số tại điểm có hoành độ $x_0=-1$ là $$f'(-1)=\dfrac{-5}{\left(2\cdot (-1)-3\right)^2}=\dfrac{-1}{5}.$$}
\end{ex}
\begin{ex}%Câu 5.%[1D7H2-3]
	Viết phương trình tiếp tuyến của đồ thị hàm số $y=x^4-4x^2+5$ tại điểm có hoành độ $x=-1$. 
	\choice
	{$y=4x-6$}
	{$y=4x+2$}
	{\True $y=4x+6$}
	{$y=4x-2$}
	\loigiai{
		Ta có $y'=4x^3-8x$, $y'(-1)=4$.\\
		Điểm thuộc đồ thị đã cho có hoành độ $x=-1$ là $M(-1;2)$.\\
		Vậy phương trình tiếp tuyến của đồ thị hàm số tại $M(-1;2)$ là
		$$y=y'(-1)(x+1)+2\Leftrightarrow y=4(x+1)+2\Leftrightarrow y=4x+6.$$}
\end{ex}
\begin{ex}%Câu 6.%[1D7H2-8]
	Cho chuyển động được xác định bởi phương trình $s=2t^3+6t^2-t$, trong đó $t$ được tính bằng giây và $s$ được tính bằng mét. Vận tốc tức thời của chuyển động tại thời điểm $t=3$ s là 
	\choice
	{\True $89$ m/s}
	{$105$ m/s}
	{$48$ m/s}
	{$20$ m/s}
	\loigiai{
		Ta có $v(t)=S'(t)=6t^2+12t-1$.\\
		Vận tốc tức thời của chuyển động khi $t=3$ s là $v(3)=6\cdot 3^2+12\cdot 3-1=89$ m/s.}
\end{ex}
\begin{ex}%Câu 7.%[1D7H2-1]
	Tính đạo hàm của hàm số $f(x)=\mathrm{e}^{2x-3}$. 
	\choice
	{\True $f'(x)=2\cdot\mathrm{e}^{2x-3}$}
	{$f'(x)=\mathrm{e}^{2x-3}$}
	{$f'(x)=-2\cdot\mathrm{e}^{2x-3}$}
	{$f'(x)=2\cdot\mathrm{e}^{x-3}$}
	\loigiai{
		Ta có $f'(x)=(2x-3)'\mathrm{e}^{2x-3}=2\cdot\mathrm{e}^{2x-3}$.}
\end{ex}
\begin{ex}%Câu 8.%[1D7H2-1]
	Hàm số $y=\ln(2x+1)$ có đạo hàm là
	\choice
	{$y'=\dfrac{2}{x\ln(2x+1)}$}
	{$y'=\dfrac{1}{2x+1}$}
	{\True $y'=\dfrac{2}{2x+1}$}
	{$y'=\dfrac{1}{(2x+1)\ln 2}$}
	\loigiai{
		Hàm số $y=\ln(2x+1)$ có đạo hàm là $y'=\dfrac{2}{2x+1}$.}
\end{ex}
\begin{ex}%Câu 9.%[1D7N3-1]
	Đạo hàm cấp hai của hàm số $y=x^6-4x^3+2x+2022$ với $x\in\mathbb{R}$ là
	\choice
	{$y''=30x^4-24x+2$}
	{\True $y''=30x^4-24x$}
	{$y''=6x^5-12x^2+2$}
	{$y''=6x^5-12x^2$}
	\loigiai{
		Ta có $y'=6x^5-12x^2+2$.
		Suy ra $y''=30x^4-24x$.}
\end{ex}
\begin{ex}%Câu 10.%[1D7N2-4]
	Cho hàm số $y=x^3-2x+1$ có đồ thị $(C)$. Hệ số góc $k$ của tiếp tuyến với $(C)$ tại điểm có hoàng độ bằng $1$ bằng
	\choice
	{$k=-5$}
	{$k=10$}
	{$k=25$}
	{\True $k=1$}
	\loigiai{
		Ta có $y'=3x^2-2$.\\
		Hệ số góc $k$ của tiếp tuyến với $(C)$ tại điểm có hoàng độ bằng $1$ bằng $k=y'(1)=1$.}
\end{ex}
\begin{ex}%Câu 11.%[1D7H2-1]
	Đạo hàm của hàm số $y=(2x-1)\sqrt{x^2+x}$ là
	\choice
	{\True $y'=\dfrac{8x^2+4x-1}{2\sqrt{x^2+x}}$}
	{$y'=\dfrac{8x^2+4x+1}{2\sqrt{x^2+x}}$}
	{$y'=\dfrac{4x+1}{2\sqrt{x^2+x}}$}
	{$y'=\dfrac{6x^2+2x-1}{2\sqrt{x^2+x}}$}
	\loigiai{
		Ta có $y'=2\sqrt{x^2+x}+\dfrac{(2x-1)(2x+1)}{2\sqrt{x^2+x}}
		=\dfrac{4x^2+4x+4x^2-1}{2\sqrt{x^2+x}}=\dfrac{8x^2+4x-1}{2\sqrt{x^2+x}}$.\\
		Vậy $y'=\dfrac{8x^2+4x-1}{2\sqrt{x^2+x}}$.}
\end{ex}
\begin{ex}%Câu 12.%[1D7N1-1]
	Cho hàm số $y=f(x)$ có đạo hàm tại điểm $x_0$. Tìm khẳng định đúng trong các khẳng định sau
	\choice
	{\True $f'(x_0)=\lim\limits_{x\to x_0}\dfrac{f(x)-f(x_0)}{x-x_0}$}
	{$f'(x_0)=\lim\limits_{x\to x_0}\dfrac{f(x)+f(x_0)}{x-x_0}$}
	{$f'(x_0)=\lim\limits_{x\to x_0}\dfrac{f(x)-f(x_0)}{x+x_0}$}
	{$f'(x_0)=\lim\limits_{x\to x_0}\dfrac{f(x)+f(x_0)}{x+x_0}$}
	\loigiai{
		Theo định nghĩa đạo hàm ta có $f'(x_0)=\lim\limits_{x\to x_0}\dfrac{f(x)-f(x_0)}{x-x_0}$.}
\end{ex}
\Closesolutionfile{ans}
\indapan{6}{ans/ans-OTC9-TN}

\cauds
\Opensolutionfile{ans}[ans/ans-OTC9-DS]
\setcounter{ex}{0}
\begin{ex}%Câu 1.%[1D7V2-4]
	Cho hàm số $y=f(x)=\dfrac{3x+2}{x+1}$. Các mệnh đề sau đúng hay sai?
	\choiceTF
	{\True $f'(x)>0$, $\forall x\neq-1$}
	{$f'(2)+f'(-3)=\dfrac{5}{36}$}
	{Tích tất cả các nghiệm của phương trình $f'(x)=25$ là $\dfrac{26}{25}$}
	{\True Tiếp tuyến của đồ thị hàm số tại điểm $M(-2;4)$ đi qua điểm $E(-5;1)$}
	\loigiai{
		\begin{itemchoice}
		\itemch \textbf{Đúng}.\\
		$f'(x)=\dfrac{(3x+2)'\cdot (x+1)-(3x+2)\cdot (x+1)'}{(x+1)^2}=\dfrac{1}{(x+1)^2}>0$, $\forall x\neq-1$.
		\itemch \textbf{Sai}.\\
		Ta có $f'(2)+f'(-3)=\dfrac{1}{(2+1)^2}+\dfrac{1}{(-3+1)^2}=\dfrac{13}{36}$.
		\itemch \textbf{Sai}.\\
		Ta có $f'(x)=25\Leftrightarrow\dfrac{1}{(x+1)^2}=25\Leftrightarrow(x+1)^2=\dfrac{1}{25}\Leftrightarrow\hoac{&x_1=-\dfrac{6}{5}\\&x_2=-\dfrac{4}{5}.}$ \\
		Nên $x_1\cdot x_2=\dfrac{24}{25}$.
		\itemch  \textbf{Đúng}.\\
		Ta có $f'(-2)=\dfrac{1}{(-2+1)^2}=1$.\\
		Phương trình tiếp tuyến của đồ thị hàm số tại điểm $M(-2;4)$ là $$y-4=1(x+2)\Leftrightarrow y=x+6.$$
		Vì $1=-5+6$ suy ra tiếp tuyến đi qua điểm $E(-5;1)$.
	\end{itemchoice}}
\end{ex}
\begin{ex}%Câu 2.%[1D7V2-8]
	Một vật chuyển động trên đường thẳng được xác định bởi công thức $s(t)=t^3-t^2+5t+2$, trong đó $t$ là thời gian tính bằng giây và $s$ là quãng đường chuyển động của vật tính bằng mét. Các mệnh đề sau đúng hay sai?
	\choiceTF
	{Vận tốc của vật tại thời điểm $t=2$ là $9$ (m/s)}
	{\True Gia tốc của vật tại thời điểm $t=3$ là $16$ (m/s$^2$)}
	{\True Quãng đường chuyển động của vật tại thời điểm mà vận tốc của vật bằng $45$ (m/s) là 70(m)}
	{\True Vận tốc của vật tại thời điểm mà gia tốc của vật bằng $34$ (m/s$^2$) là $101$ (m/s)}
	\loigiai{
		\begin{itemchoice}
		\itemch  \textbf{Sai}.\\
		Vận tốc của vật tại thời điểm $t$ là $v(t)=s'(t)=3t^2-2t+5$.\\
		Vận tốc của vật tại thời điểm $t=2$ là $v(2)=3\cdot 2^2-2\cdot 2+5=13$ (m/s).
		\itemch  \textbf{Đúng}.\\
		Gia tốc của vật tại thời điểm $t$ là $a(t)=s''(t)=6t-2$.\\
		Vận tốc của vật tại thời điểm $t=3$ là $a(3)=6\cdot 3-2=16$ (m/s$^2$).
		\itemch  \textbf{Đúng}.\\
		Vận tốc của vật bằng $45$ (m/s) tại thời điểm $t$ suy ra $$v(t)=45\Leftrightarrow 3t^2-2t+5=45\Leftrightarrow\hoac{&t=4 \,(\text{thỏa mãn})\\&t=-\dfrac{10}{3}\quad(\text{loại}).}$$
		Quãng đường chuyển động của vật tại thời điểm $t=4$ là\\ $s(4)=4^3-4^2+5\cdot 4+2=70$ (m).
		\itemch  \textbf{Đúng}.\\
		Gia tốc của vật bằng $34$ (m/s$^2$) tại thời điểm $t$ suy ra $$a(t)=34\Leftrightarrow 6t-2=34\Leftrightarrow t=6.$$
		Vận tốc của vật tại thời điểm $t=6$ là $v(6)=3\cdot 6^2-2\cdot 6+5=101$ (m/s).
	\end{itemchoice}}
\end{ex}
\begin{ex}%Câu 3.%[1D7V2-5]
	Cho đồ thị $(C)$: $y=\dfrac{2x+3}{x-2}$; $M(1;-5)$; điểm $N\in(C)$ có hoành độ bằng $3$. Khi đó các mệnh đề sau đúng hay sai?
	\choiceTF
	{\True Hệ số góc của tiếp tuyến với $(C)$ tại điểm $M$ bằng $-7$}
	{Tung độ tiếp điểm $N$ bằng $-5$}
	{\True Tiếp tuyến với $(C)$ tại $M$ song song với đường thẳng $d_1\colon y=-7x+17$}
	{\True Có hai tiếp tuyến với $(C)$ vuông góc với đường thẳng $d_2\colon 16x-7y+35=0$}
	\loigiai{
		\begin{itemchoice}
		\itemch  \textbf{Đúng}.
		Ta có $y'=\dfrac{-7}{(x-2)^2}\Rightarrow$ hệ số góc của tiếp tuyến với $(C)$ tại điểm $M$ là $$k=y'(1)=-7.$$
		\itemch  \textbf{Sai}.
		$x_N=3\Rightarrow y_N=9$.
		\itemch  \textbf{Đúng}.
		Vì tiếp tuyến $\Delta\colon y=-7x+2$ với $(C)$ tại $M$ và $d_1\colon y=-7x+17$ có hệ số góc bằng nhau nên $\Delta\parallel d_1$.
		\itemch  \textbf{Đúng}.
		Gọi $(x_0; y_0)$ là tọa độ của tiếp điểm.\\
		Vì tiếp tuyến với $(C)$ vuông góc với đường thẳng $d_2\colon 16x-7y+35=0$ nên $$f'(x_0)\cdot\dfrac{16}{7}=-1\Leftrightarrow\dfrac{-7}{(x_0-2)^2}\cdot\dfrac{16}{7}=-1\Leftrightarrow(x_0-2)^2=16\Leftrightarrow\hoac{&x_0=6\\&x_0=-2.}$$ 
		Vậy có hai tiếp tuyến cần tìm.
	\end{itemchoice}}
\end{ex}
\begin{ex}%Câu 4.%[1D7V2-8]
	Một vật chuyển động trên đường thẳng được xác định bởi công thức $s(t)=t^3-3t^2+7t-2$, trong đó $t>0$ và tính bằng giây và $s$ là quãng đường chuyển động được của vật trong $t$ giây tính bằng mét. Khi đó
	\choiceTF
	{\True Tốc độ của vật tại thời điểm $t=2$ là $7$ (m/s)}
	{\True Gia tốc của vật tại thời điểm $t=2$ là $6$ (m/s$^2$)}
	{Gia tốc của vật tại thời điểm mà vận tốc của chuyển động bằng $16$ (m/s) là $10$ (m/s$^2$)}
	{\True Thời điểm $t=1$ (s), tại đó vận tốc của chuyển động đạt giá trị nhỏ nhất}
	\loigiai{
		Ta có $s'(t)=3t^2-6t+7$ và $s''(t)=6t-6$.
		\begin{itemchoice}
		\itemch  \textbf{Đúng}.\\
		Vận tốc của vật tại thời điểm $t=2$ là $v(2)=s'(2)=3\cdot 2^2-6\cdot 2+7=7$ (m/s).
		\itemch  \textbf{Đúng}.\\
		Gia tốc của vật tại thời điểm $t=2$ là $a(2)=v'(2)=s''(2)=6\cdot 2-6=6$ (m/s$^2$).
		\itemch  \textbf{Sai}.\\
		Vận tốc của chuyển động bằng $16$ (m/s$^2$) tại thời điểm $t$ nghĩa là
		$$v(t)=s'(t)=16\Leftrightarrow 3t^2-6t+7=16\Leftrightarrow\hoac{&t=3 (\text{thỏa mãn})\\&t=-1 (\text{loại}).}$$
		Gia tốc của vật tại thời điểm $t=3$ là $a(3)=v'(3)=s''(3)=6\cdot 3-6=12$ (m/s$^2$).
		\itemch  \textbf{Đúng}.\\
		Vận tốc của chuyển động có phương trình $v(t)=3t^2-6t+7$ là một parabol, có đỉnh là
		$S\left(-\dfrac{b}{2a};-\dfrac{\Delta}{4a}\right)\Rightarrow S(1;4)$ và hệ số $a=3>0$ nên hàm số có giá trị nhỏ nhất bằng 4 tại $t=1$
		(hoặc sử dụng máy tính Casio để tìm min).\\
		Vậy tại thời điểm $t=1$ thì vận tốc của chuyển động đạt giá trị nhỏ nhất bằng $4$ (m/s).
	\end{itemchoice}}
\end{ex}
\Closesolutionfile{ans}
\indapan{3}{ans/ans-OTC9-DS}

\Opensolutionfile{ans}[ans/ans-OTC9-KQ]
\caukq
\setcounter{ex}{0}
\begin{ex}%Câu 1.%[1D7V2-1]%[2D1V3-6]%[2D1N3-6]
	Cho hàm số $y=\dfrac{x^2+4x-1}{2x+3}$ có đạo hàm $y'=\dfrac{ax^2+bx+c}{(2x+3)^2}$. Khi đó giá trị của biểu thức $P=a^2+b^2+c^2$ bằng
	\shortans{$236$}
	\loigiai{\allowdisplaybreaks\begin{eqnarray*}
		y'&=&\dfrac{\left(x^2+4x-1\right)'\cdot (2x+3)-(2x+3)'\cdot\left(x^2+4x-1\right)}{(2x+3)^2}\\
		&=&\dfrac{(2x+4)(2x+3)-2\left(x^2+4x-1\right)}{(2x+3)^2}\\
		&=&\dfrac{4x^2+14x+12-2x^2-8x+2}{(2x+3)^2}\\
		&=&\dfrac{2x^2+6x+14}{(2x+3)^2}.
		\end{eqnarray*}
		$ \Rightarrow a=2$; $b=6$; $c=14$.\\
		Vậy $P=a^2+b^2+c^2=236$.}
\end{ex}
\begin{ex}%Câu 2.%[1D7V2-1]
	Cho hàm số $y=f(x)=\sqrt{\tan x+\cot x}$. Tính $f'\left(\dfrac{\pi}{4}\right)$.
	\shortans{$0$}
	\loigiai{
		\textbf{Cách 1:}\allowdisplaybreaks\begin{eqnarray*} f'(x)&=&\dfrac{1+\tan^2x-\left(1+\cot^2x\right)}{2\sqrt{\tan x+\cot x}}\\
			&=&\dfrac{\tan^2x-\cot^2x}{2\sqrt{\tan x+\cot x}}.
			\end{eqnarray*}
		Suy ra $f'(0)=\dfrac{\tan^2\dfrac{\pi}{4}-\cot^2\dfrac{\pi}{4}}{2\sqrt{\tan\dfrac{\pi}{4}+\cot\dfrac{\pi}{4}}}=0$.\\
		\textbf{Cách 2: }Sử dụng máy tính Casio.}
\end{ex}
\begin{ex}%Câu 3.%[1D7V2-8]
	Cho biết điện lượng truyền trong dây dẫn theo thời gian biểu thị bởi hàm số $Q(t)=2t^2+t$, trong đó $t$ được tính bằng giây và $Q$ được tính theo Culông. Tính cường độ dòng điện tại thời điểm $t=4$ (s).
	\shortans{$17$}
	\loigiai{
		Ta có $Q'(t)=\left(2t^2+t\right)'=4t+1$ nên cường độ dòng điện tại thời điểm $t=4$(s) là $Q'(4)=4\cdot 4+1=17$ (C/s).}
\end{ex}
\begin{ex}%Câu 4.%[1D7V2-8]
	Cho chuyển động thẳng xác định bởi phương trình $s(t)=-t^3+3t^2+9t$, trong đó $t$ tính bằng giây và $s$ tính bằng mét. Vận tốc của chuyển động tại thời điểm gia tốc triệt tiêu bằng bao nhiêu?
	\shortans{$12$}
	\loigiai{
		Vận tốc của chuyển động chính là đạo hàm cấp một của quãng đường $$v(t)=s'(t)=-3t^2+6t+9.$$
		Gia tốc của chuyển động chính là đạo hàm cấp hai của quãng đường $$a(t)=s''(t)=-6t+6.$$
		Gia tốc triệt tiêu khi $s''(t)=0\Leftrightarrow t=1$.\\
		Khi đó vận tốc của chuyển động là $s'(1)=12$ m/ s.}
\end{ex}
\begin{ex}%Câu 5.%[1D7V2-8]
	Một vật chuyển động thẳng được xác định bởi phương trình $s(t)=t^3-2t^2+2t$, trong đó $t$ tính bằng giây và $s$ tính bằng mét. Tính gia tốc của vật tại thời điểm mà vận tốc của vật bằng $17$ (m/s).
	\shortans{$14$}
	\loigiai{
		Ta có vận tốc của vật tại thời điểm $t$ là $v(t)=s'(t)=3t^2-4t+2$.\\
		Gia tốc của vật tại thời điểm $t$ là $a(t)=s''(t)=6t-4$.\\
		Vận tốc của vật bằng $17$ (m/s) tại thời điểm $t$ suy ra $$v(t)=17\Leftrightarrow 3t^2-4t+2=17\Leftrightarrow\hoac{&t=3 (\text{thỏa mãn})\\&t=-\dfrac{5}{3}\quad(\text{loại}).}$$
		Gia tốc của vật tại thời điểm $t=3$ là $a(3)=6\cdot 3-4=14$ (m/s$^2$).}
\end{ex}
\begin{ex}%Câu 6.%[1D7V2-8]
	Một vật chuyển động thẳng xác định bởi phương trình $s(t)=-\dfrac{1}{3}t^3+3t^2+7t-2$, trong đó $t$ tính bằng giây và $s$ tính bằng mét. Vận tốc của vật đạt giá trị lớn nhất bằng bao nhiêu?
	\shortans{$16$}
	\loigiai{
		Ta có vận tốc của vật tại thời điểm $t$ là $v(t)=s'(t)=-t^2+6t+7$.\\
		Khi đó ta có $v(t)=-(t-3)^2+16\leq 16$ suy ra hàm số $v(t)$ đạt giá trị lớn nhất bằng $16$ khi và chỉ khi $t-3=0\Leftrightarrow t=3$.\\
		Vậy vận tốc của vật đạt giá trị lớn nhất bằng $16$ tại $t=3$.}
\end{ex}
\Closesolutionfile{ans}
\indapan{6}{ans/ans-OTC9-KQ}